\section{The service receipt on the selected public run}\label{ex-centered-public-service}

The finite-volume moment estimate uses one service certificate at each period, reference scale, and terminal volume. The public route builds these certificates from a single funded run and its same-history multiscale record. This section follows that finite producer and spells out how its fields are used; it does not assume that a new history can be selected for each volume.

\begin{theorem}[Same-history service certificates]\label{thm:ex-selected-run-service}
Fix the selected run for a simple gauge group and a campaign album on its trail, with its one family and gauge doors. For each grade $d$, let $k_d$ denote its selected reference scale. Assume the funded record on this trail supplies (i) a single transition history with retention clauses at every level and a terminal endowment; (ii) the scale room and bank budgets, the initial positivity, exponential scale, radius-splitting, and support-size inequalities of Proposition~\ref{prop:ex-geyser-service-fields}, and a nondecreasing occupancy schedule $\rho_i$; (iii) for each grade $d$, period $p$, and indices $k_d\leq k\leq N$, a height bound $X=X(d,k,N,p)\geq1$, the coefficient and activity-price inequalities and adjacent weight-growth inequalities at levels at least $k$, and the terminal bank, footprint-cover, and local-section estimates stated below; (iv) for each grade $d$, period $p$ and terminal volume $N\geq k_d$, a calibrated finite music whose voices are linear (the coordinate, regular and forced parts of each voice depend complex-linearly on the coefficient, and its marked support does not depend on the coefficient), adjacent door estimates, and its terminal forced-profile identity; and (v) gradewise thresholds, positive scale prices with adjacent growth, and finite local tone tickets on those same doors. Then for every period $p$, reference level $k\geq1$, and terminal volume $N$, the certificate has retention and terminal endowment on the restricted history; room, bank, all four stated scale inequalities, and monotone occupancy at levels $i\geq k$; and, for every grade with $k_d\leq k\leq N$, the calibrated finite transfer and the height, coefficient, activity, adjacent-weight, terminal-bank, footprint-cover (at most one), and uniform local-section fields specified in the hypotheses. More explicitly, for each such grade and every $i\geq k$, writing $\Gamma_i$ for the columns, $\ell_i$ for the block side, $\zeta_i=\zeta_d(i)$ and $a_d$ for the orbit cost, the stock witness satisfies
\[
 \zeta_i\Gamma_i\leq1,\quad \psi_i\geq0,\quad
 2a_d(\psi_i+2\Gamma_{i+1})\zeta_{i+1}\leq1,\quad
 \phi_i\leq1,\quad \phi_i\ell_i^d\Gamma_{i+1}\zeta_i\leq\zeta_{i+1}.
\]
At $N$, let $P_N$ be the terminal budget data, $E_N$ its support exponent, $\theta_N$ its positive threshold, and put $\operatorname{fullCost}_d(P_N,X)=X\theta_N^{-\deg F_d}$. The witness also satisfies
\[
 X\geq1,\quad e^{E_N+2}\operatorname{fullCost}_d(P_N,X)\zeta_N\leq u_N,\quad
 C_{B_N}(\operatorname{footprint}_d)\leq1,\quad
 \|\operatorname{places}_d(\operatorname{axis}_j,0,U)\|\leq X
\]
for every frame index $j$ and terminal configuration $U$. The calibrated transfer uses the same doors and terminal profile identity as in the hypotheses. The lower restriction $k\geq1$ is retained from the public service constructor. All fields refer to the same selected history. The height witness can depend on all four displayed indices; its local-section bound is uniform over the finite frame indices and configurations at that volume.
\end{theorem}

\begin{proof}
The selected run fixes one trail $T$, one family $C$, transfer stories, and gauge doors at every adjacent scale; the purchase record supplies the multiscale ledgers. The public route applies the global Geyser service producer to exactly this tuple; its per-period Pond is a derived, weaker interface. Its output contains a retained history, the tab schedule, a terminal keeper, marked portfolios, purchases, and weak specifications. The history's cost field contains the numerical room and bank estimates, the scale inequalities, monotone occupancy, gradewise height data, and terminal music records. All are attached to the same selected run.

Fix $p,k,N$. The service certificate is assembled from those records as follows.

First, restrict the retained history to levels $i\geq k$. At each adjacent scale in this range, the selected spring and purchased tariff give a local comparison. The endpoint support exponents are at least $100$, so the finite interpolation constructor returns both its local gate ticket and the duty estimate. These are exactly the clauses needed for the retained-history field; the terminal endowment is the terminal clause of the same history.

Next, use the room and bank estimates from the same clock. Restrict those initial positivity, exponential scale, radius-splitting, and support-size inequalities to indices $i\geq k$; their bounds are unchanged. The occupancy schedule remains monotone under this restriction.

Finally, when $k\leq N$, for every grade $d$ with reference level at most $k$, the height record at $(d,k,N,p)$ supplies a constant $X$ and the data needed downstream: the coefficient and activity-price bounds and adjacent weight-growth inequalities at every level at least $k$, the terminal bank bound, the terminal footprint-cover estimate, and the uniform bound on the local-section insertions. The abstract certificate allows $X$ to depend on all four indices. In the selected construction, the uniform insertion estimate gives a stronger witness $X_d$ depending only on the grade: it bounds every frame-axis insertion at every finite volume, period and configuration, while the fitted stencil gives the required footprint cover. The terminal music record supplies the other part: its finite music begins with the calibrated factory entries, has this linearity at the grade reference, pays each adjacent transition through the same doors, and satisfies the terminal forced-profile identity. Joining these witnesses gives, for each such grade, the complete terminal service record. When $k>N$, the stock clause has no obligation. The same construction applies at every allowed $(p,k,N)$ and introduces no new history; the other service fields retain their stated index dependence.

The same-history per-label construction applies this service record to a finite list of observables. The affordability induction in Section~\ref{ex-centered-construction} propagates each label's initial payment through the adjacent coefficient and scale-price bounds. It uses the prescribed transported weights and one music on the selected history; it does not choose a new history at each scale.

For the moment estimate, take a finite set $V$ of labels of one fixed grade $d$, with positive marker weights $w_i$. The per-label construction supplies compatible marked activities; finite backward propagation of their terminal product then gives a served marked representation at the reference scale. Let $X(U)$ be its represented marked polynomial, $A$ the local Wilson representative, $\sigma=\operatorname{scalar}(A)$ its common scalar, and $\mathcal T_k$ the vacuum cost per site. Write $\operatorname{toll}_w(I)=\prod_{i\in I}w_i$ and $\operatorname{price}_w(Y)=\sum_I|Y_I|\operatorname{toll}_w(I)$ for its mode toll and polynomial price. Suppose the terminal data satisfy the hypotheses of the selected root-test estimate: $k<N$, the local section is sound, the representation is selected, weights are positive, the column-weighted prices obey $w_i\Gamma_k\,\mathrm{paid}_i\leq\zeta_d(k)$, and the label sites have the stated migration, occupancy, and separation bounds. The strict inequality $k<N$ identifies the propagated support exponent at level $k$ with the current support exponent; the service certificate itself is also defined when $k=N$. Here $\Gamma_k\geq1$ is the column count and each $\mathrm{paid}_i\geq0$, so the same inequality also implies the unweighted initial payment $w_i\,\mathrm{paid}_i\leq\zeta_d(k)$ used in the preceding induction. If $a=\zeta_d(k)=\mathfrak f_k\leq1$ is the selected label-price allowance, the column-weighted bound gives the vault budgets $x=3a$ and $y=a$ by the finite response estimate. The local representative is lawful, has nonnegative coupling and pointwise nonnegative evaluation, and its support exponent is nonnegative. At this selected level put $s=\widehat h_k=\widehat\beta_k^{-1/2}$. The selected delayed-coupling bound gives $\widehat\beta_k>1$, hence $0<s\leq1$; the selected vacuum threshold is $s$ times the exponential of a square, so $s\leq\theta_k$. This choice uses the same selected history and is the radius in its treasury receipt. Its served representation has reserve $f=1$ and integrable mode coefficients, as furnished by the finite transition receipt. The pointwise price estimate is
\[
 \operatorname{price}_w(X(U))\leq
 \exp\!\left(\operatorname{scalar}(A)+|\Lambda_k|\bigl(\operatorname{exCost}(A)+\kappa\bigr)
                 +|V|(x+y+2)\right).
\]
Indeed this is the pointwise price bound \eqref{ex-centered-steer} of Section~\ref{ex-centered-construction}, with $s=|\Lambda_k|$ sites, $x=3a$, $y=a$ and reserve $f=1$, combined there with the lawfulness bound $E(U)\leq\operatorname{scalar}(A)+|\Lambda_k|\operatorname{exCost}(A)$; the two extra units are charged by that reserve. The denominator lower bound is
\[
 \int \operatorname{eval}(A)\,d\mu
 \geq \exp\!\left(\operatorname{scalar}(A)-|\Lambda_k|
       \bigl(\operatorname{exCost}(A)+\beta\,\operatorname{baseBudget}
                     +4d_G\log(16/s)\bigr)\right),
\]
where $\beta\geq0$ and $d_G$ is the real dimension of the ambient matrix representation space. The same vacuum treasury receipt bounds the per-site bill
\[
 2\operatorname{exCost}(A)+\kappa+\beta\,\operatorname{baseBudget}
       +4d_G\log(16/s)\leq\mathcal T_k.
\]
The selected Wilson transition identities preserve the integral of the
unmarked density, so the denominator here equals the terminal Wilson
partition function, denoted $Z_{N,p}$. The common scalar cancels. Integrating a fixed mode
coefficient, using its positive toll, and dividing by this denominator gives
the finite local bound
\[
\begin{split}
\frac{\left|\int X_I\,d\mu\right|}{\int \operatorname{eval}(A)\,d\mu}
       \operatorname{toll}_w(I)
 &\leq \exp\!\Bigl\{|\Lambda_k|\bigl(2\operatorname{exCost}(A)+\kappa
          +\beta\,\operatorname{baseBudget}\\
 &\hspace{47mm}+4d_G\log(16/s)\bigr)+|V|(4a+2)\Bigr\}.
\end{split}
\]
This follows directly from the displayed numerator and denominator: $|X_I|\operatorname{toll}_w(I)\leq\operatorname{price}_w(X)$ pointwise, so $\left|\int X_I\right|\operatorname{toll}_w(I)\leq\int\operatorname{price}_w(X)$; the prior is normalized, and $x+y+2=4a+2\leq6$. The finite marked-coefficient estimate is therefore at most $\exp(|\Lambda_k|\mathcal T_k+6|V|)$. In the public root test, let $M$ be the number of cells at root scale $m$, with one replicated label per cell and common weight $w>0$, and put $q=(a_m/a_k)^4$, the four-dimensional volume ratio. The centered spacer estimate applies to this moment bound.
The aligned side-count identity is $|\Lambda_k|=Mq$. Thus
the finite-volume exponent is $M(q\mathcal T_k+6)$. Taking the normalized
$M$th root and dividing by $w$ gives the centered root-test estimate
\[
                       \exp(q\mathcal T_k+6)/w.
\]
This is the marked-copy spacer estimate~\eqref{ex-centered-opened-root} of Section~\ref{ex-centered-construction}.
The treasury contribution survives with factor $q$ in the root exponent; it
is not canceled by the finite-volume denominator. Finite backward Wilson
integration preserves each mode integral, as proved in
Section~\ref{ex-selected-birth-receipt}.

The service record, the same-history propagation and the finite local price all come from the one funded record assumed in (i)--(v); the displayed computations show how its clauses enter the bound, and no unpriced per-volume assumption is added.
\end{proof}

The construction of the funded record itself is not proved in this chapter. It is input I-C02-funded-record of the package ledger; its ingredients are the selected campaign purchase, the transition tariff, the delayed scale choices, the same-trail transfer data, the climbing doors and the finite-transfer records.
